%! Function Model Construction and Application — Unit 1, Lesson 1.8 — Homework (Answer Key)
\documentclass[10pt]{article}
\usepackage{precalculus-article}
\usepackage{precalculus-key}   % pulls in -boxes; do NOT also load -boxes
\newcommand{\TallMath}[1]{$\displaystyle #1\rule[-1.4em]{0pt}{3.2em}$}

\begin{document}

\pageheader{Unit 1, Lesson 1.8}{Homework --- Answer Key}
\namedateperiod

\vspace{0.1in}

\begin{enumerate}[itemsep=18pt]

\item \textbf{Model from Restrictions.} A farmer has 120 feet of fencing to enclose a
      rectangular garden. Both pairs of opposite sides require fencing (no wall).
      Let $w$ = width in feet.

      \begin{enumerate}[label=(\alph*), itemsep=8pt]
        \item Write the constraint equation and solve for length $\ell$ in terms of $w$.

              \vspace{0.06in}
              Constraint: \ans{$2w + 2\ell = 120$} \quad $\Rightarrow$ \quad $\ell = $ \ans{$60 - w$}

        \item Write the area function $A(w)$.

              \vspace{0.06in}
              $A(w) = $ \ans{$w(60 - w) = -w^2 + 60w$}

        \item State the domain restriction.

              \vspace{0.06in}
              Domain: \ans{$0 < w < 60$}

        \item Find the width and length that maximize area.

              \vspace{0.06in}
              $w_{\max} = $ \ans{30} ft \quad $\ell_{\max} = $ \ans{30} ft \quad $A_{\max} = $ \ans{900} ft$^2$
      \end{enumerate}

\begin{teachernote}
Problem 1: the unconstrained garden (no wall) produces a square at maximum area. This contrasts with Lesson 1.8's barn example where the optimal is not square. Good discussion point.
\end{teachernote}

\item \textbf{Rational Inverse-Proportion Model.} The gravitational force between two objects
      is modeled by $F(d) = \dfrac{k}{d^2}$, where $d$ is the distance in meters and $k = 4\times10^{14}$
      N$\cdot$m$^2$.

      \begin{enumerate}[label=(\alph*), itemsep=8pt]
        \item State the domain restriction and describe what happens to $F$ as $d \to \infty$.

              \vspace{0.06in}
              Domain: \ans{$d > 0$} \quad As $d \to \infty$, $F(d) \to $ \ans{$0$}.

        \item Compute $F(2 \times 10^7)$ and $F(4 \times 10^7)$. What is the ratio $F(2\times10^7)\,/\,F(4\times10^7)$?

              \vspace{0.06in}
              $F(2\times10^7) = \dfrac{4\times10^{14}}{4\times10^{14}} = $ \ans{$1$ N} \quad $F(4\times10^7) = \dfrac{4\times10^{14}}{16\times10^{14}} = $ \ans{$0.25$ N}

              Ratio: \ans{$4:1$} \quad This illustrates the \ans{inverse-square} law.
      \end{enumerate}

\item \textbf{Apply the Model --- Average Rate of Change.} A ball's position (in feet) after
      $t$ seconds is modeled by $s(t) = -16t^2 + 64t + 5$.

      \begin{enumerate}[label=(\alph*), itemsep=8pt]
        \item Evaluate $s(1)$ and $s(3)$.

              \vspace{0.06in}
              $s(1) = -16 + 64 + 5 = $ \ans{53} ft \quad $s(3) = -144 + 192 + 5 = $ \ans{53} ft

        \item Compute the average rate of change of $s$ on $[1, 3]$. Include units.

              \vspace{0.06in}
              $\text{ARC} = \dfrac{s(3)-s(1)}{3-1} = \dfrac{{\color{keyred}\mathbf{0}}}{{\color{keyred}\mathbf{2}}} = {}$ \ans{$0$}

              Units: \ans{feet per second (ft/s)}

        \item Interpret the ARC: what does it tell you about the ball's motion?

              \vspace{0.06in}
              \ansline{An ARC of 0 ft/s means the ball is at the same height at $t=1$ and $t=3$ s; the ball went up and came back down symmetrically over this interval (the vertex is at $t=2$ s).}
      \end{enumerate}

\item \textbf{Interpret with Units.} A model gives the volume of water $V(t)$ in gallons
      in a tank $t$ minutes after draining begins. The average rate of change on $[0, 10]$ is
      $-15$ gal/min.

      \begin{enumerate}[label=(\alph*), itemsep=8pt]
        \item Is the tank filling or draining? How do you know from the sign of the ARC?

              \vspace{0.06in}
              \ansline{The tank is draining. The ARC is negative, meaning volume decreases on average as time increases.}

        \item If $V(0) = 200$ gal, estimate $V(10)$ using the average rate of change.
              Is this an exact value or an approximation? Why?

              \vspace{0.06in}
              $V(10) \approx $ \ans{50} gal. This is an \ans{approximation} because \ansline{the ARC is an average; the actual rate of draining may not be constant, so $V(10)$ could differ.}
      \end{enumerate}

\item \textbf{Piecewise Model.} A downtown parking garage charges:
      \begin{itemize}[itemsep=2pt]
        \item A flat fee of \$3.00 for the first hour (or any part thereof).
        \item \$2.00 per hour for each additional hour after the first.
      \end{itemize}

      Let $C(t)$ = total cost in dollars for $t$ hours of parking, $t > 0$.

      \begin{enumerate}[label=(\alph*), itemsep=8pt]
        \item Write the piecewise model. (Treat hours as a continuous variable.)

              \[
                C(t) = \begin{cases}
                  {\color{keyred}\mathbf{3}} & \text{if } 0 < t \leq 1 \\[6pt]
                  {\color{keyred}\mathbf{3 + 2(t-1)}} & \text{if } t > 1
                \end{cases}
              \]

        \item Compute $C(0.5)$, $C(1)$, $C(2.5)$, and $C(4)$.

              \vspace{0.06in}
              $C(0.5) = $ \ans{\$3.00} \quad $C(1) = $ \ans{\$3.00} \quad $C(2.5) = $ \ans{\$6.00} \quad $C(4) = $ \ans{\$9.00}

        \item Is $C(t)$ a polynomial function? Explain.

              \vspace{0.06in}
              \ansline{No; $C(t)$ is a piecewise function defined by different rules on different intervals, not a single polynomial expression valid for all $t > 0$.}
      \end{enumerate}

\item \textbf{AP-Style Multiple Choice.} A function model for a company's revenue (in thousands
      of dollars) is $R(x) = -x^2 + 8x$, where $x$ is the number of units sold (in hundreds),
      $0 \leq x \leq 8$. What is the average rate of change of $R$ on $[2, 6]$?

      \medskip
      \begin{enumerate}[label=(\Alph*), itemsep=3pt]
        \item $0$ \textcolor{keyred}{\textbf{$\leftarrow$ correct}} \qquad \textbf{(B)} $4$ \qquad \textbf{(C)} $-4$ \qquad \textbf{(D)} $8$
      \end{enumerate}

      \vspace{0.06in}
      Answer: \ans{A} \quad Show work:

      \vspace{0.08in}
      \ansline{$R(2) = -4+16=12$; $R(6)=-36+48=12$; ARC $= (12-12)/(6-2) = 0$. The vertex is at $x=4$, so the function is symmetric about $x=4$ on $[2,6]$.}

\end{enumerate}

\begin{reflectionbox}
What strategy do you use first when deciding whether a model should be polynomial or rational?
\ansline{Check whether the context describes a relationship where one quantity grows by powers (polynomial) or decreases as the other increases according to a reciprocal power (rational/inverse proportion).}
\end{reflectionbox}

\end{document}
